%% This LaTeX-file was created by <imt> Wed May  6 17:11:32 1998
%% LyX 0.12 (C) 1995-1998 by Matthias Ettrich and the LyX Team

%% Do not edit this file unless you know what you are doing.
\documentclass{article}
\usepackage[T1]{fontenc}
\usepackage{a4wide}
\usepackage{graphics}

\makeatletter


%%%%%%%%%%%%%%%%%%%%%%%%%%%%%% LyX specific LaTeX commands.
\newcommand{\LyX}{L\kern-.1667em\lower.25em\hbox{Y}\kern-.125emX\spacefactor1000}

%%%%%%%%%%%%%%%%%%%%%%%%%%%%%% Textclass specific LaTeX commands.
\newenvironment{lyxlist}[1]
  {\begin{list}{}
    {\settowidth{\labelwidth}{#1}
     \setlength{\leftmargin}{\labelwidth}
     \addtolength{\leftmargin}{\labelsep}
     \renewcommand{\makelabel}[1]{##1 \hfill}}}
  {\end{list}}

\makeatother

\begin{document}


\title{PROJECT REQUIREMENT SPECIFICATION}


\author{System Design Group Seven}

\maketitle
\vspace{0.3cm}
{\centering \includegraphics{../web/images/EUcrest-official-noborder.ps} \par}
\vspace{0.3cm}


\newpage
\section{Product Task}

Ding Dong Doorentry Systems Ltd have commissioned our group to develop an idea
for a new product. The product will be a door entry system that will allow a
person within a building to use their own PC to view and communicate to the
person at the door. Our task is to provide a prototype unit and the accompanying
software for this project. The Market Potential for this device can be broken
down down as:


\subsection{Residential }

Information from AmeriTech (a CCTV security firm) showed that you are 5 times
less likely to be burgled if you have some form of surveillance scheme. The
scheme we aim to provide would have to meet the competeing needs of specific
user groups. They can be broken down:


\subsubsection{Blocks of Flats }

The increase of PC's in the home may allow this product to become an alternative
to the standard door-entry system. According to a representative of Castle Security(a
local security company) it would be likely to be unwanted by technophobes and
those on a limited income for the forseeable future as at the moment the majority
of people are quite happy with the system they have; they can already identify
anyone who comes to their door. A secondary issue would be some people in a
block wanting the technology while others may not. Would the system work in
tandem with present technology.


\subsubsection{Accommodation with a servitor}

The servitors office would contain the PC and during the day people could be
let in by sight only. During the night a form of security could be set on the
main enterances. The potential here is for accommodation like halls of residence,
the security could be contained on a Dallas Ring which would hold enough information
to also act as a matriculation card.


\subsubsection{Stand Alone Housing }

Unlikely to be a market in smaller property, people will usually answer their
door or look out through a window to see who it is. There may be a market in
countries like the USA where there is a significant gun culture, and residents
already have a fear of intruders. The main market at present would be properties
standing in a large area ie. estate houses. They can have entrances that are
set a significant distance away from the property and could use the improved
security of a cisual communication. They are also likely to have the need and
the money to invest in this form of security


\subsection{Commercial}

This will be our proposed target market, in light of the Cullen Inquiry into
the Dunblane report And the 1995 NHS Executive reccommendation the need for
low cost security has significantly increased. Lord Cullen stated

\begin{quotation}
It is regognised that there should be arrangements for security to prevent unauthorised
access to the school.
\end{quotation}
The Health and Safety at Work Regulations(1974) and as amended (1992) requires 

\begin{quotation}
That every employer shall make a suitable and sufficient assesment of ... the
risks to the health and safety of persons not in his employment arising out
of or in connection with the conduct by him or his undertaking. 
\end{quotation}
Employers have a legal obligation to ensure the security of their staff.


\subsubsection{Education}

By encouraging schools to place this system at the entry points to their site
visitor movement can be controlled. During specific times the doors would be
required to be left open, eg 8.15 - 9.15, to allow pupil and teacher access
but outside of these times access to the site will be controlled inside the
school. Allowing the system to be configurable will allow this policy to be
changed ie. the children could each have their own dallas rings. Outside of
school times the site will be closed to all but those with security access,
which should have the added effect of cutting down vandalism. In situations
were people are using part of the school out of hours then they will have to
possess a security pass.


\subsection{NHS }

The system could be integrated to the front doors of wards that demand security
for example maternity wards. Visitors will only be able to enter these areas
if they have a specific reason to be there. The Problem with existing NHS security
include

\begin{itemize}
\item security being lax because no-one monitors the vide camers.
\item security is bypasses by people who pretend to be some one else
\end{itemize}
The system we aim to supply must be quick and easy to use, since there may be
emergency situations where the time is criticall, particuly a cardiac arrest
alarm. The demands on the system will involve Doctors needing instant access
to wards whilst still keeping a level of security that will satisfy parents.


\newpage
\section{Minimum System}


\subsection{System Description}

Any visitor to the site will be able to request access by ringing a door bell.
If the user is at their PC they can run the software to use the camera. The
software will display one still image pictue that will be regularly updated,
this picture will be save to file. The user will be able to communicate with
the visitor from the PC. If the user is satisfied that this person should access
to the site then they will be able to send a signal to open the door.

When there is no reply from the user the visitor will be enable access to the
site if they know the correct pin number.


\subsection{Software description}

The Client-side features:

\begin{itemize}
\item Will display a picture which will regularly be updated
\item The picture will be saved to a file
\item Will handle duplex audio
\item The database information will be stored in a text format
\item Will supply a way to open the door
\item Ther will be an operational mode control (deactivate automatic system, indicate
a person is leaving the PC)
\end{itemize}
The general software wil include 

\begin{itemize}
\item Keep alive indicators
\end{itemize}

\subsection{Hardware features}

The unit will 

\begin{itemize}
\item Have a keypad for inputing pin numbers
\item Handle duplex audio with no compression
\item There will be a lock relay
\end{itemize}

\newpage
\section{Intended System}


\subsection{System Description}

Any visitor to the site will be able to request access by ringing a door bell.
The program will pop up, giving the user instant communication with the visitor.
The software will display one still image picture that will be regularly updated,
this picture will be saved to a file. The user will be able to communicate with
the visitor from the PC. If the user is satisfied that this person should access
to the site then they will be able to send a signal to open the door. The user
will also be able to acccess diagnostic information at any time so that they
can find out the current state of the system.

When there is no reply from the user the visitor will be enable access if they
are able to place the correct Dallas ring in the unit and know the correc correct
pin number to go with it. The pins could be optional.


\subsection{Software description}

The Client-side features

\begin{itemize}
\item Will include database management software
\item The program will pop up when the doorbell rings.
\item There will be a stored masterkey sequence contained in the non-volatile part
of memory.
\end{itemize}
The general software wil include 

\begin{itemize}
\item A report on the diagnostics of the system.
\end{itemize}

\subsection{Hardware features}

The unit will comprise of a:

\begin{itemize}
\item Dallas ring reader
\item A door sensor
\end{itemize}

\newpage
\section{Ideal System}


\subsection{System Description }

The PC will display a continuous stream of video. The user will be able to stop
the frame at any point, this will allow the user to only see the most regognisable
image the camera provides. The database information will contain a picture of
the person whose pin number has been typed allowing the user to see if the picture
on record is the same as the person attempting to use the pin.If the picture
is the same the visitor must have authorised access if its different appropriate
steps can be taken. The user will then be able to communicate with the visitor
and enable access to the site. 

When there is no user at the PC, requesting entry will enable the PC to reply
with a recorded message of the users choosing. The camera will allso store some
frames to keep a record of who has accessed the site. Access can then be gained
by a Dallas ring placed in a reader and typing in the correct access code. The
PC will contain a database of information (pin numbers) with everything running
through netscape. The system will contain a modem to allow a remote site to
control the doors.

Therefore a large site can be controlled by a LAN linking all of cameras and
centrally controlled and viewed through the internet at one specific site.


\subsection{Software description}

The Client-side features:

\begin{itemize}
\item Will show an official picture based on the users ID
\item We will implement an http server
\end{itemize}
The general software wil include 

\begin{itemize}
\item The PC will determine wether the user is absent
\end{itemize}

\subsection{Hardware features}

The unit will comprise of a:

\begin{itemize}
\item Will have some form of audio hardware compression.
\end{itemize}

\newpage
\section{Hazards + Threats}

\begin{lyxlist}{00.00.0000}
\item [\textbf{Fault:}]Failure of power system
\end{lyxlist}
\textbf{Reported by:} 

\begin{lyxlist}{00.00.0000}
\item [\textbf{Likelihood:}]High in a remote area, unlikely in urban areas.
\item [\textbf{Consequences:}]Undesirable, the system will shut down and will need
to be booted up manually
\end{lyxlist}
\textbf{Response to failure:}

The doors can be opened from the inside by the mechanical lock. From the outside
by a key.

\begin{lyxlist}{00.00.0000}
\item [\textbf{Fault:}]Failure of electonic entry system
\end{lyxlist}
\textbf{Reported by:}

\begin{lyxlist}{00.00.0000}
\item [\textbf{Likelihood:}]low
\item [\textbf{Consequences:}]Entry will only take place manually.
\item [\textbf{Response}]\textbf{to failure:}
\end{lyxlist}
The door will open with a manual key and an alarm will go off. If the person
entering is genuine they will know how to overide the alarm. If it is a break-in
then the alarm will act as a signalling for help.

\end{document}
